\documentclass[11pt]{article}

\usepackage[margin=1in]{geometry}
\usepackage{parskip}
\usepackage{booktabs}
\usepackage{array}
\usepackage{hyperref}
\usepackage{xcolor}
\usepackage{enumitem}

\title{Progress Report --- Reliability of LLMs for Vulnerability Detection}
\author{}
\date{}

\begin{document}

\begin{center}
{\LARGE \textbf{Progress Report --- Reliability of LLMs for Vulnerability Detection}}\\[1em]
\begin{tabular}{ll}
\textbf{Date:} & [To be written] \\
\textbf{Prepared for:} & [Professor name] \\
\end{tabular}
\end{center}

\vspace{1em}
\hrule
\vspace{1em}

\begin{abstract}
\noindent This report covers my LLM vulnerability-detection project on the CWE-415 Juliet corpus.
Section~\ref{sec:background} is the background: what the study is trying to do, what the corpus
is, and the terms I use throughout. The progress starts in Section~\ref{sec:bottomline}.

\noindent The short version is this. GPT-4o gets 53.8\% accuracy on this corpus. The usual
explanation is label leakage --- names like \texttt{bad()} and \texttt{badSink} give the answer
away. The evidence does not settle that cleanly. On the 1,306 files that declare no
\texttt{bad*()} function, the model still answers ``vulnerable'' 89.6\% of the time. Paired
discrimination across the corpus is 11.4\%, below the 25\% you would get from coin flips. But a
stricter check does show the cue carries some weight. On the 276 files where no
\texttt{bad}-prefixed token survives anywhere, accuracy roughly triples, to 0.304. So the fair
claim is that the model badly under-uses the cue, not that it ignores it. On top of that, I found
two coverage gaps in the anonymization pipeline. They leave 68\% of the corpus only partly
anonymized. Until those are closed, the leakage-versus-reasoning question stays open. This week
is about closing them, re-checking the leak rate, and finishing a short list of low-cost fixes
and experiments.
\end{abstract}

\vspace{0.5em}

\section{Project Background}
\label{sec:background}

This section is background only. It restates what the study is trying to do, for anyone who has
not seen the manuscript draft. The progress itself starts in Section~\ref{sec:bottomline}.

\subsection*{The study}

Large Language Models (LLMs) have recently been proposed as tools for automated
vulnerability detection; however, there is still limited evidence that they can consistently
identify security-critical bugs. In this work, we attempt to present a comprehensive empirical study on the
reliability of LLMs for vulnerability detection. Specifically, initially we intend to examine how small,
syntactically neutral code transformations --- such as renaming API or function identifiers,
removing comments, or altering whitespace and formatting --- affect model performance. Using a
synthetic subset of the Juliet Test Suite, we would like to evaluate four representative LLM
families: GPT-3.5-Turbo, GPT-4o-Mini, Claude, and Llama-3. We would further compare the
effectiveness of zero-shot, few-shot, and chain-of-thought (CoT) prompting strategies.

The central question is whether an LLM's apparent detection ability survives transformations
that change nothing about the program's semantics. If accuracy collapses when identifiers are
renamed or comments are stripped, the model was reading surface cues rather than reasoning about
memory behaviour --- which matters, because a real codebase does not label its own bugs.

\subsection*{The corpus}

Experiments use the \textbf{Juliet Test Suite}, NIST's synthetic C/C++ benchmark. Juliet is
convenient because ground truth is exact and the classes are balanced, but it has a property
that drives most of this report. Nothing in Juliet is written by hand. Every testcase is
generated from a template, and the generator writes the answer into the names it produces. The
vulnerable function is literally called \texttt{bad()}. Its safe twin is called \texttt{good()}.
The construct around them usually carries the whole CWE identifier
(\texttt{CWE415\_Double\_Free\_\_malloc\_free\_char\_81}). So the ground truth is not something a
model has to work out from the memory behaviour. It is spelled out in the identifiers, before the
model reads a single statement. Each testcase is emitted as an
\texttt{\_omitbad}/\texttt{\_omitgood} pair differing by only a few lines. The work reported here
uses the CWE-415 (double free) and CWE-416 (use after free) subsets; the wider corpus statistics
are in Table~\ref{tab:corpus}.

\begin{table}[h!]
\centering
\begin{tabular}{lrrrr}
\toprule
CWE & Total files & Good files & Vulnerable files & Header files \\
\midrule
CWE-416 & 1,040 & 520 & 520 & 0 \\
CWE-415 & 3,368 & 1,644 & 1,644 & 80 \\
CWE-191 & 9,577 & 4,722 & 4,722 & 133 \\
CWE-190 & 13,108 & 6,444 & 6,444 & 220 \\
CWE-122 & 19,747 & 9,630 & 9,630 & 487 \\
CWE-121 & 16,049 & 7,871 & 7,871 & 307 \\
\bottomrule
\end{tabular}
\caption{Unprocessed dataset statistics, Juliet Test Suite.}
\label{tab:corpus}
\end{table}

\subsection*{The two experimental variables}

\subsubsection*{Cleaning setups}
The cleaning setups control how much non-semantic material the model sees. This is the main
variable the study is built around:

\begin{itemize}[leftmargin=*]
  \item \texttt{clean\_setup\_0} --- raw dataset, unmodified. All comments, macros, and
  formatting preserved. This is the baseline, and it is the only setup that has been run so far,
  which is why every number in Section~\ref{sec:data} comes from it.
  \item \texttt{clean\_setup\_1} --- moderate cleaning: remove comments, rename all user-defined
  functions and namespaces, remove macros except \texttt{\#include}, collapse blank lines, drop
  empty or invalid source files.
  \item \texttt{clean\_setup\_2} --- aggressive normalization extending setup 1: newlines
  collapsed to single spaces, tabs replaced with spaces, minimizing superficial code variation.
\end{itemize}

The renaming in setups 1 and 2 is the \emph{anonymization pipeline} --- the Joern/libclang stage
that rewrites Juliet's identifiers into neutral names. A function named \texttt{bad()} gives the
answer away. A model can answer ``vulnerable'' without reading the code at all. It is a common assumption that renaming removes
that shortcut. So every claim in this report about leakage depends on the renaming being
complete.

\subsubsection*{Prompting strategies}
The prompting strategy is the second variable. There are three:

\begin{itemize}[leftmargin=*]
  \item zero-shot,
  \item few-shot, and
  \item chain-of-thought (CoT).
\end{itemize}

Each one comes in two variants: a \emph{basic} variant that asks ``is this code vulnerable?'',
and a \emph{CWE-specific} variant that asks ``is this code vulnerable to CWE-415?''. Every result
reported below uses \texttt{zero\_shot\_cwe} --- zero-shot prompting with the CWE-specific
question.


\subsection*{Where this report fits}

The plan is straightforward. First, anonymize the corpus. Then run each model on every cleaning
setup and prompting strategy. If accuracy falls once the cues are gone, that drop tells us how
much the model was leaning on them.

The rest of this report is about why I am not there yet. The baseline results were stranger than
I expected. And when I checked the anonymization step, I found two gaps in what it actually
renames. Those gaps have to be fixed first. Until they are, comparing one cleaning setup against
another does not tell us much.

\section{Bottom Line}
\label{sec:bottomline}

I started from the usual explanation. GPT-4o does badly on the CWE-415 corpus: 53.8\% accuracy,
and it answers ``vulnerable'' on 90.4\% of the files. The assumption is that this is label
leakage. The model sees a name like \texttt{bad()}, \texttt{badSink}, or
\texttt{CWE415\_Double\_Free\_\_...} and reads the answer off it.

The numbers do not really support that. There are 1,306 files with no \texttt{bad*()} function in
them at all, although most of them still have the word ``bad'' somewhere, left behind in the
header comments. GPT-4o calls 89.6\% of them vulnerable. So I narrowed it down further. In the
276 files where no \texttt{bad}-prefixed token survives anywhere, accuracy roughly triples, to
0.304.

So the model is not ignoring the cue. It is using it far less than it should, given how clean a
signal it is. I'm still working through what that means. Part of the reason is two gaps I found
in how the anonymization pipeline decides what counts as an identifier to rename:

\begin{itemize}[leftmargin=*]
  \item \textbf{Gap 1 --- the rule picks symbols by name, and Juliet's C++ testcases slip past
  it.} The pipeline renames a symbol only if its name matches the regex \texttt{.*CWE.*}
  (\texttt{src/joern/joern\_symbols.sc}, lines 8--9). Juliet's C-style testcases put the full
  \texttt{CWE415\_Double\_Free\_\_...} prefix on the function name. Those match, and they get
  renamed. The C++-style testcases do it differently. They wrap the same logic in a namespace
  and drop the prefix from everything inside it, so the names are just \texttt{bad()},
  \texttt{badSink}, \texttt{badGlobal}. Juliet's own source comment says this is deliberate ---
  the namespace already makes the name unique. The regex never sees any of them. The result:
  2,288 of 3,368 files (68\%) still contain at least one label-bearing identifier after
  ``anonymization.''

  \item \textbf{Gap 2 --- class and struct names are never looked at.} The collection script
  queries \texttt{cpg.method}, \texttt{cpg.local}, and \texttt{cpg.namespaceBlock}. It never
  queries \texttt{cpg.typeDecl}. So a name like
  \texttt{class CWE415\_Double\_Free\_\_malloc\_free\_char\_81\_bad} is invisible to it. That
  name matches the \texttt{.*CWE.*} pattern perfectly. It just never becomes a candidate,
  because nothing collects it in the first place. This affects 80 files.
\end{itemize}

What this means in practice: I can only trust the ``the model under-uses leaked names'' finding
on the part of the corpus where anonymization actually worked. That is why both gaps have to be
closed this week. It holds no matter which way the analysis eventually points.

\section{What the Data Shows}
\label{sec:data}

All of this comes from the \texttt{zero\_shot\_cwe} runs already in
\texttt{results/all\_results.csv}. One caveat first. Every row there is \texttt{clean\_setup\_0}.
No cleaning was applied at all, so the model was reading prompts with the CWE label sitting right
there in the code. Even with the answer handed to it, GPT-4o did not do well.

\begin{table}[h!]
\centering
\begin{tabular}{lrrrrrrrrr}
\toprule
Model / CWE & Total & Acc & Prec & Rec & F1 & TP & TN & FP & FN \\
\midrule
gpt-4o / CWE-415 & 3,288 & 53.8\% & 52.1\% & 94.2\% & 0.6707 & 1,548 & 220 & 1,424 & 96 \\
gpt-4o / CWE-416 & 1,040 & 52.0\% & 51.0\% & 99.8\% & 0.675 & 519 & 22 & 498 & 1 \\
\bottomrule
\end{tabular}
\caption{Published results, \texttt{zero\_shot\_cwe} strategy.}
\end{table}

The F1 scores look almost respectable. The confusion matrices do not. On CWE-415 there are 1,424
false positives against only 220 true negatives. On CWE-416 it is 498 false positives, 22 true
negatives, and a single false negative. The model said \textit{vulnerable} on 90.4\% of CWE-415
samples and 97.8\% of CWE-416 samples. An F1 near 0.67 on a balanced task is what you get from a
model that says yes to almost everything.

That is the aggregate picture. It gets clearer when I break CWE-415 down by what the file
actually contains:

\begin{table}[h!]
\centering
\begin{tabular}{lrrr}
\toprule
Subset & $n$ & Predicted ``vulnerable'' & Accuracy \\
\midrule
File declares \texttt{bad*()} & 1,294 & 100.0\% & 1.000 \\
File declares only \texttt{good*()} & 1,306 & 89.6\% & 0.104 \\
\quad --- of which, no \texttt{bad}-prefixed token at all & 276 & 69.6\% & 0.304 \\
\bottomrule
\end{tabular}
\caption{Behavioural breakdown, GPT-4o on CWE-415.}
\end{table}

The first row is not evidence of much. The model predicts vulnerable on almost everything, so any
group that is 100\% actually-vulnerable scores 1.000 for free. The second row is closer to a real
test, but not quite. These files declare no \texttt{bad*()} function, but most of them still say
``bad'' somewhere in the header comments Juliet leaves behind (\texttt{BadSource:},
\texttt{BadSink:}, and so on). The model calls them vulnerable 89.6\% of the time anyway. The
third row is the one that matters --- the 276 files with no \texttt{bad}-prefixed token anywhere.
There the positive rate drops and accuracy roughly triples, to 0.304. So the model does use the
cue. It just uses it far less than it should, given that the cue separates the two classes almost
perfectly.

Aggregate metrics can hide a pattern like this. So I also computed \textbf{paired
discrimination}. Juliet generates each testcase as an \texttt{\_omitbad}/\texttt{\_omitgood} pair
that differs by only a few lines. I scored a pair as correct only when the model got both halves
right.

\begin{center}
PairedDiscrimination $= 187 / 1{,}644 = \textbf{11.4\%}$
\end{center}

Independent coin flips would land at 25\%. Scoring \textit{below} chance is not noise. It is what
you see when a model gives the same answer no matter what it is shown.

The false positives make this concrete. In several cases the model names the safe function
correctly in its own explanation, and still answers \textit{vulnerable}:

\begin{quote}
\texttt{..\_new\_delete\_array\_char\_42\_omitbad.cpp} $\rightarrow$ ``The function
\texttt{goodB2GSource} deletes ...''\\
\texttt{..\_malloc\_free\_int64\_t\_74b\_omitbad.cpp} $\rightarrow$ ``The function potentially
frees memory ...''\\
\texttt{..\_new\_delete\_class\_05\_omitbad.cpp} $\rightarrow$ ``The code contains a double free
vulnerabi...''
\end{quote}

It finds the right function and still gets the label wrong. That is not a lexical shortcut
failing. It looks more like a default answer the model falls back to, whatever it just read.

\section{Why This Forces a Code Change This Week}
\label{sec:fix}

The two gaps in Section~\ref{sec:bottomline} have the same root cause. In the pipeline, what gets
anonymized is decided by a regex. Right now that regex is \texttt{.*CWE.*}, so a symbol is
renamed only if its \textit{name} contains ``CWE'' somewhere. I think that approach has a
problem. It means guessing every naming convention Juliet might use, in every CWE directory,
forever. I
could widen the regex to catch more variants --- \texttt{cwe}, \texttt{bad}, \texttt{good},
\texttt{sink}, \texttt{source}, \texttt{flaw}. That would catch more of them, but it is still a
guess. And I would have to audit it again for the next corpus.

The fix I want to make this week is simpler than a better regex. Stop matching on the name. Match
on where the symbol is declared instead. Every function, variable, namespace, and type declared
in a file inside the input directory gets renamed, with \texttt{main} as the only exception.
Under that rule there is nothing left to miss. It does not matter what Juliet named
\texttt{bad()}, because the rule never looks at the name.

It also generalizes better. In CWE-121/122 the leaking identifiers are \texttt{dataBadBuffer} and
\texttt{dataGoodBuffer}. There is no ``CWE'' substring anywhere in them. A location-based rule
catches those on day one. A name-based rule catches them only after someone notices.

Concretely, the change touches two files:

\begin{table}[h!]
\centering
\begin{tabular}{>{\raggedright\arraybackslash}p{5.5cm}>{\raggedright\arraybackslash}p{9.2cm}}
\toprule
File & Change \\
\midrule
\texttt{joern\_symbols.sc} & Replace the name-regex filter with a declared-in-corpus filter;
add \texttt{cpg.typeDecl} collection so class/struct names are seen at all. \\
\addlinespace
\texttt{03\_apply\_\allowbreak renames\_\allowbreak libclang.py} & Add class/struct declaration
and type-reference kinds so the newly-collected symbols are actually rewritten, not just
detected. \\
\bottomrule
\end{tabular}
\caption{Change set for the anonymization fix.}
\end{table}

One consequence I want to flag to you directly. This makes the corpus much more aggressively
anonymized than before --- \texttt{func\_000412} calling \texttt{func\_000413}, with none of
Juliet's original naming left. If part of the goal is to mirror how the original benchmark
authors preprocessed their data, that is a real divergence, not a side effect. I would rather
state it in the methods section than let a reviewer find it. I think it is the right call for
measuring detection ability, but I want your read before I treat it as settled.

On urgency: until this is fixed, I cannot honestly write ``every user-declared identifier is
replaced'' in a methods section. It is not true right now, and it is the kind of claim a reviewer
checks in five minutes. That holds no matter how the leakage-versus-reasoning question resolves,
which is why it comes before everything else on my list this week.

\section{Plan for the Upcoming Week}

Roughly in the order I plan to do them:

\begin{enumerate}[leftmargin=*]
  \item \textbf{Fix the anonymization pipeline.} The two-file change from
  Section~\ref{sec:fix}: the declaration-site filter and \texttt{cpg.typeDecl} collection in
  \texttt{joern\_symbols.sc}, plus the matching rewrite support in
  \texttt{03\_apply\_renames\_libclang.py}. This is the main thing this week is for.

  \item \textbf{Re-run the leak audit on the fixed pipeline.} Right now 68\% of files still leak
  the label after anonymization. I want to see that number drop to zero before I trust anything
  downstream of it.

  \item \textbf{Fix the \texttt{'Prediction'} key bug in \texttt{evaluator.py}.} A ten-minute
  fix. It is silently mislabeling the precision column, and I would rather catch it now than
  after I regenerate the result tables.

  \item \textbf{Add paired discrimination as a reported metric.} I compute it by hand at the
  moment. Putting it in \texttt{evaluator.py} means every future run reports it automatically,
  including the runs on the fixed pipeline.

  \item \textbf{Run \texttt{clean\_setup\_1} and \texttt{clean\_setup\_2}.} Both already exist
  and neither has ever been run. This is pure inference time, no new code. I am curious what
  stripping the \texttt{POTENTIAL FLAW} comment does to precision. That comment is currently
  \textit{anti}-correlated with the vulnerable label (Section~\ref{sec:data}), so the direction
  of the effect is a genuine open question, not a foregone conclusion.
\end{enumerate}

If the pipeline fix lands early enough in the week, I would also like to re-run the GPT-4o
evaluation on the newly anonymized corpus, and see whether the 89.6\% false-positive rate on
clean files holds up. That is a stretch goal, not a commitment. Items 1--3 come first.

\section{Open Question for Discussion}

There is one thing I would like your input on. The fixed pipeline anonymizes much more
aggressively than the current one. Every declared identifier gets replaced, not just the ones
with ``CWE'' in the name. I have not tracked down how other papers using Juliet handle this, so I
have no baseline for whether that level of anonymization is standard or unusual. I would rather
flag the gap than assume it does not matter. Is it worth chasing down before I finalize the
pipeline, or is it fine to state as a deliberate choice either way?

\end{document}
